\Def Кольцо -- набор из множества «чисел» и операций сложения и
умножения, для которых верны все классические свойства этих
операций: ассоциативность, коммутативность,
дистрибутивность\dots Однако нет обратимости по умножению (нельзя
делить).

\Def Кольцо $\Z_m$ -- структура над числами $0, 1, \dots , m - 1$, где все
операции ведутся по модулю $m$

\Note Это наивное определение, которое в общем некорректно, но дает понимание того, с чем мы работаем.

\Def Отрицательное число в $\Z_m$: \space
$-b\,=\,
\begin{cases}
0,\:b\:=\:0\\
m - b,\:\text{иначе}
\end{cases}$

\Def $a^{-1} = b$ в $\Z_m$ тогда и только тогда, когда $a\,\cdot\,b\,\equiv\,1\,(mod \,m)$.

\Def Поле — набор из множества «чисел» и операций сложения и умножения, для которых верны все классические свойства этих операций: ассоциативность, коммутативность, дистрибутивность и обратимость.

\Note Это наивное определение, которое в общем некорректно, но дает понимание того, с чем мы работаем.

\Thbd (Малая теорема Ферма). Пусть $m$ -- простое, тогда
$1 = a^{m - 1}$ в $\Z_m$.

\Consequence Если $p$ -- простое, то $\Z_p$ -- поле, обратное тоже верно. То
есть для каждого ненулевого $a \in \Z_p$ существует $a^{-1}$

\Thbd  (Эйлер) Пусть $m \in \N$. Тогда для $a \in \Z_m$ существует  $a^{-1}$ тогда и только тогда, когда $gcd(a, m) = 1$. Более того, если $gcd(a,\,m)\,=\,1$, то $a^{-1} = a^{\varphi(m)-1}$ в $\Z_m$.


\pagebreak